% part: intuitionistic-logic
% chapter: introduction
% section: syntax

\documentclass[../../../include/open-logic-section]{subfiles}

\begin{document}

\olfileid{int}{int}{syn}

\olsection{د شهودي منطق نحو}

د شهودي منطق نحو د بياني منطق د نحو په شان ده۔ په کلاسيک بياني
منطق کښې ځينې نښلوونکي د نورو په وسيله تعريفېدے شي؛ د بېلګې په توګه،
$!A \lif !B$ د $\lnot !A \lor !B$ په وسيله، يا $!A \lor !B$ د
$\lnot(\lnot !A \land \lnot !B)$ په وسيله تعريفولے شو۔ نو د کلاسيک
منطق په ځينو وړاندې‌کونو کښې ځينې نښلوونکي د دغو تعريفونو لنډې نښې
وي۔ په شهودي منطق کښې داسې نۀ ده، مګر دوه استثناوې شته: $\lnot !A$
د $!A \lif \lfalse$ لنډه نښه کېدے شي، او ډېر ځله همداسې تعريفېږي۔
خو بيا، البته، $\lfalse$ پخپله بايد تعريف شوې نښه نۀ وي!{} همدارنګه،
$!A \liff !B$ د کلاسيک منطق په شان د $(!A \lif !B) \land (!B \lif
!A)$ په توګه تعريفېدے شي۔

د بياني شهودي منطق !!^{formula}s له \emph{!!{propositional variable}s}،
د قضيوي ثابت~$\lfalse$، او \emph{منطقي نښلوونکو} څخه جوړېږي۔ موږ لرو:

\begin{enumerate}
\item د !!{propositional variable}s يو !!^a{denumerable} سټ~$\PVar$:
  $\Obj p_0$، $\Obj p_1$، \dots
  \item د !!{falsity} قضيوي ثابت~$\lfalse$۔
\item منطقي نښلوونکي: $\land$ (عطف)، $\lor$ (فصل)، $\lif$ (شرطي)۔
\item د ليک نښې: پرانيستونکې قوس «(»، تړونکې قوس «)»، او کامه۔
\end{enumerate}

\begin{defn}[فارمول]
\ollabel{defn:formulas} د بياني شهودي منطق د \emph{!!{formula}s}
سټ~$\Frm[L_0]$ په استقرايي ډول داسې تعريفېږي:
\begin{enumerate}
\item $\lfalse$ يو اټومي !!{formula} دے۔

\item هر !!{propositional variable}~$\Obj p_i$ يو اټومي
  !!{formula} دے۔

\item که $!A$ او $!B$ !!{formula}s وي، نو $(!A \land
  !B)$ يو !!a{formula} دے۔

\item که $!A$ او $!B$ !!{formula}s وي، نو $(!A \lor !B)$
  يو !!a{formula} دے۔

\item که $!A$ او $!B$ !!{formula}s وي، نو $(!A \lif !B)$
  يو !!a{formula} دے۔

\tagitem{limitClause}{له دې پرته بل هېڅ څۀ !!a{formula} نۀ دے۔}{}
\end{enumerate}
\end{defn}

د پورته لومړنيو نښلوونکو تر څنګ، لاندې \emph{تعريف شوې} نښې هم
کاروو: $\lnot$ (نفي) او $\liff$ (!!{biconditional})۔ په دې تعريف
شويو عاملونو جوړ شوي فارمولونه داسې لوستل کېږي:
\begin{enumerate}
\item $\lnot !A$ د $!A \lif \lfalse$ لنډه نښه ده۔

\item $!A \liff !B$ د $(!A \lif !B) \land (!B
  \lif !A)$ لنډه نښه ده۔
\end{enumerate}

که څۀ هم $\lnot$ په رسمي ډول لنډه نښه ګڼل کېږي، کله کله به په
تعريفونو کښې د~$\lnot$ لپاره ښکاره قاعدې او مادې داسې ورکوو لکه
چې لومړنۍ نښه وي۔ دا زياتره د تمريني مسئلو د بيان لپاره کوو۔

\end{document}
